\ifdefined\iscombined
	\lecturetitle{Finite State Processes}
	\setcounter{section}{2}
\else
\documentclass{article}
\usepackage{amsthm}
\usepackage{comment}
\usepackage{amsmath}
\usepackage{amsfonts}
\usepackage{sectsty}
\usepackage{hyperref}
\usepackage{fancyhdr}
\usepackage[top=1in,bottom=1in,left=1in,right=1in]{geometry}
\usepackage{float}
\usepackage{graphicx}
\usepackage{tabularx}
\usepackage{mathtools}

%TODO: Consider changing formatting so examples/notes are not side by side

\title{Lecture 2 - Finite State Processes}
\author{Matthew Farah, \href{mailto:farahm11@mcmaster.ca}{$\langle$farahm11@mcmaster.ca$\rangle$}, 400468251}
\date{\today}

\hypersetup{
        colorlinks=true,
        linkcolor=blue,
        filecolor=magenta,
        urlcolor=blue,
        citecolor=blue,
	anchorcolor=blue
}

\fancyhead{}
\fancyhead[L]{Matthew Farah, \href{mailto:farahm11@mcmaster.ca}{$\langle$farahm11@mcmaster.ca$\rangle$}, 400468251}
\fancyhead[R]{Lecture 2 - Finite State Processes}

\newenvironment{example}[1][]{
	\vspace{1em}
	\large{\par\noindent\textbf{Ex.}}
}{\vspace{2.5em}}

\newenvironment{solution}[1][]{
	\par\noindent\textbf{\large{Solution:}}
	\par
	\vspace{1.5em}
}{\vspace{.5em}}

\renewcommand{\thesubsection}{\thesection.\alph{subsection}}
\renewcommand{\thesubsubsection}{\thesection.\alph{subsection}.\Roman{subsubsection}}
\makeatletter

\sectionfont{\fontsize{12}{12}\bfseries}
\subsectionfont{\fontsize{10}{12}\normalfont}
\subsubsectionfont{\fontsize{10}{12}\normalfont}

\begin{document}

\maketitle
\pagestyle{fancy}

\fi

\begin{itemize}
	\item FSPs are algebraic models used to represent the behaviour of sequential and concurrent processes.
	\begin{itemize}
		\item FSPs are composed of \texttt{PROCESSES} and \texttt{actions}.
		\begin{itemize}
			\item An \texttt{action} is a thing which the system does.
			\begin{itemize}
				\item An \texttt{action} must be atomic (can only do one thing).
			\end{itemize}
			\item A \texttt{PROCESS} perform an action and calls another process.
		\end{itemize}
	\end{itemize}
\end{itemize}

\begin{example}
	Represent a process which performs action \texttt{x} and then calls another processes \texttt{P} as an FSP.
\end{example}

\begin{solution}
	\begin{align*}
		(\mathtt{x} \rightarrow \mathtt{P})
	\end{align*}
\end{solution}

\begin{itemize}
	\item Repetitive behaviour is modelled by recursively calling processes.
\end{itemize}

\begin{example}
	Represent a simple on-off switch as an FSP.
\end{example}

\begin{solution}
	\begin{align*}
		\mathtt{SWITCH} &= \mathtt{OFF} \\
		\mathtt{OFF} &= \left( \mathtt{on} \rightarrow \mathtt{OFF} \right) \\
		\mathtt{ON} &= \left( \mathtt{off} \rightarrow \mathtt{OFF} \right) \\
	\end{align*}
\end{solution}

\begin{itemize}
	\item \emph{Process substitution} is the act of replacing processes with their definitions to produce more succinct FSPs
\end{itemize}

\begin{example}
	Using process substitution, use the FSP for Example 2 to product an FSP for an on-off switch with only two processes.
\end{example}

\begin{solution}
	\begin{align*}
		\mathtt{SWITCH} &= \mathtt{OFF} \\
		\mathtt{OFF} &= \left( \mathtt{on} \rightarrow ( \mathtt{off} \rightarrow \mathtt{OFF}) \right) \\
	\end{align*}
\end{solution}

\begin{itemize}
	\item FSPs do not need to be deterministic.
	\item In FSP \emph{choices} allow a process to execute processes non-deterministically.
	\begin{itemize}
		\item Denoted as ``\texttt{|}''.
	\end{itemize}
\end{itemize}

\begin{example}
	Represent a process which can perform action \texttt{x} then execute process \texttt{P} or perform \texttt{y} then execute \texttt{Q}.
\end{example}

\begin{solution}
	\begin{align*}
		(\mathtt{x} \rightarrow \mathtt{P} \; | \; \mathtt{y} \rightarrow \texttt{Q})
	\end{align*}
\end{solution}

\begin{itemize}
	\item A \emph{boolean condition} is a statement which can either be true or false.
	\begin{itemize}
		\item A boolean condition is said to have a \emph{finite domain} if there exists a finite number of cases needed to be checked to evaluate it as true or false.
		\begin{itemize}
			%TODO: Not entirely sure this is correct
			\item For example, the boolean condition \texttt{x = 0} has a finite domain, however \texttt{x > 0} does not.
		\end{itemize}
	\end{itemize}
	\item A \emph{guarded action} is an action in a choice can only be executed so long as some boolean condition is met.
	\begin{itemize}
		\item The condition allowing a guarded action to be taken is specified with the \texttt{when} keyword.
		\item The boolean condition must have a finite domain.
	\end{itemize}
\end{itemize}

\begin{example}
	Represent a process which can perform action \texttt{x} then \texttt{P} \underline{ONLY IF} a condition \texttt{B} is true or perform \texttt{y} then execute \texttt{Q}.
\end{example}

\begin{solution}
	\begin{align*}
		(\mathtt{when} \; \mathtt{B} \; \mathtt{x} \rightarrow \mathtt{P} \; | \; \mathtt{y} \rightarrow \texttt{Q}) \\
	\end{align*}
\end{solution}

\begin{itemize}
	\item The \emph{alphabet} of a process is defined as the set of all possible actions that a process can perform.
	\begin{itemize}
		\item The FSP representation of a representation implicitly defines the alphabet of a process.
		\item If two FSPs have no actions in common, their alphabets are said to be \emph{disjoint}.
		\item An action is said to be \emph{shared} between two processes if it exists in both of their alphabets.
	\end{itemize}
	\item An \emph{alphabet extension} is used to specify additional actions available to a process not defined in its FSP.
\end{itemize}

\begin{itemize}
	\item If two processes are running concurrently, they are said to be running in \emph{parallel composition}.
	\item The actions of a system composed of two concurrently running processes is an interleaving of the actions of each process.
	\begin{itemize}
		\item The order of the actions taken by each system \underline{must} be preserved.
		\begin{itemize}
			\item For example, if a process \texttt{P} can execute action \texttt{a} then \texttt{b}, and process \texttt{Q} can execute \texttt{c}, a valid output of the system \texttt{(P||Q)} could be $\mathtt{a} \rightarrow \mathtt{c} \rightarrow \mathtt{b}$ but $\mathtt{b} \rightarrow \mathtt{a} \rightarrow \mathtt{c}$ would be invalid, since in \texttt{P}, \texttt{a} must be executed before \texttt{b}.
		\end{itemize}
	\end{itemize}
	\item If two processes running in parallel composition have shared actions, then the shared actions must be executed simultaneously by both processes.
	\item FSPs use the parallel composition operator to declare two processes as running concurrently.
	\begin{itemize}
		\item Denoted as ``\texttt{||}''
		\item The parallel composition operator is both associative and commutative.
	\end{itemize}
	\item Petrinets allow for much easier graphical representations of concurrent systems than LTSs.
\end{itemize}

\ifdefined\iscombined
\newpage
\else
  \end{document}
\fi
